\section{Introduction}
\label{sec:intro}

% VALUE: (1) context. Client NPUs run LLMs as compositions of per-operator designs on a spatial array, issued by one in-order CP; kernels and mappings are well served; the expert designs still leave columns idle in prefill and read KV at half rate in decode; gap: how the composition is issued has received far less attention
Neural processing units (NPUs) such as the AMD XDNA~2 in Ryzen\texttrademark{} AI processors bring large language model (LLM) inference to client machines~\citep{rico2024xdna}.
Their AI Engine array is spatial, so each operator of a model runs on its own group of columns, and the host composes the operators' work into a sequence of commands that a single in-order command processor (CP) executes~\citep{rico2024xdna,linux-amdnpu}.
Compilers and mapping frameworks already map single operators and fused groups onto the array~\citep{wang2026mlirair,hunhoff2025iron,symons2025stream}, and expert end-to-end designs exist for prefill and decode~\citep{wang2026mlirair,fastflowlm-repo}.
Yet these expert designs leave much of the array idle.
In MLIR-AIR's prefill of Llama-3.2-1B, the device is busy for over 90\% of the time-to-first-token (TTFT), but attention and the matrix multiplications (GEMMs) run on disjoint columns that take turns, and attention's share of TTFT grows from 12\% to 49\% between 512 and 4{,}096 tokens.
In decode, long-context attention reads its key-value (KV) cache at half the DRAM read rate that the platform sustains.
However, while prior work has optimized what each operator computes and where it runs, how the operators are composed and issued through the command processor has received far less attention.
% SOURCE: busy 92-97% of TTFT, attention 12.2% -> 49.1% of TTFT at 512 -> 4,096 tokens (amd-npu e3acc8a prefill-e2/decomp.txt; Section 3, Figure busy).
% SOURCE: decode attention reads KV at 28.4 GB/s = 51% of the same-batch read cap (vault lab-decode-long-context (5)).

% VALUE: (2) existing approaches and why they fall short: one operator per command with a hand-fixed split; packing into one dispatch is standard practice but on an in-order CP gives no overlap; a wait on one group holds back the other; the right split changes with the workload. Points to Figure 1.
\begin{figure}[t]
    \centering
    % VALUE: the same operators under three compositions: one per command (groups take turns), packed in one command in turn (a wait on attention holds back the GEMMs), and issued first with synchronizations in predicted completion order (ours; both groups busy)
    \figph{Schematic: CP instruction order and column activity for one attention operator (columns 0 and 7) and three GEMMs (columns 1--6): (a) one operator per command; (b) one command, operators in turn; (c) one command, synchronizations in predicted completion order (ours).}
    \caption{CP instruction order and the activity of the column groups for one attention operator and three GEMMs: (a) one operator per command, (b) one command with the operators in turn, and (c) one command with synchronizations in predicted completion order.}
    \label{fig:order}
\end{figure}
Expert designs fix the split of the array by hand and issue one operator per command, so the column groups execute one after another (Figure~\ref{fig:order}a).
Packing several operators into one dispatch is standard practice for amortizing launch overhead, as in CUDA Graphs~\citep{nvidia-cudagraphs} and XRT's run lists~\citep{xrt-repo}.
On an in-order CP, however, packing alone creates no overlap, because a synchronization on one group delays the issue of every later instruction, whichever group it feeds (Figure~\ref{fig:order}b).
The same delay occurs inside a single operator, where the existing attention stream waits on one column before it starts the other.
Furthermore, the split that suits a short prompt computed from scratch does not suit a long prompt appended to a cached context, as in an agent's turn, so one hand-fixed split leaves part of the array idle in one regime or the other.

% VALUE: (3) this work: the CP's instruction stream is the schedule; the composer makes two decisions automatically on unmodified designs (instruction order, spatial split), with legality derived from the instructions; each with its headline number; combined result last
In this work, we treat the CP's instruction stream as the schedule of the spatial array and compose it automatically from an existing design's per-operator streams, with no change to the design's source, compiler, or driver.
First, our composer issues every group's work before the CP waits on any group and places each synchronization at its predicted completion, so that disjoint column groups run concurrently (Figure~\ref{fig:order}c).
This instruction order lowers TTFT by 1.4--1.8$\times$ on two model families with bit-identical outputs, and by a steady 1.5--1.6$\times$ on prompts appended to a cached context.
Second, the composer chooses the split of columns between attention and the GEMMs for each regime, which favors two attention columns for prompts computed from scratch and four for appends to a long context.
On an append of 4{,}096 tokens to a 16{,}384-token context, the four-column split lowers the time of the composed schedule by a further 1.9$\times$ on both models.
Because the chosen splits stay resident, switching between them costs under 1\% of an append.
Third, the composer derives which operators may share a command from the instructions themselves, and it lifts the one-GEMM-per-command limit of the existing argument binding by placing many operators' buffers in shared arenas.
Together, the instruction order and the split lower the time of such an append by 3.0$\times$ on Qwen3-4B.
% SOURCE: from scratch, composed vs the existing host program on the default split: Llama-3.2-1B 1.41x at 2,048 and 1.70x at 4,096 tokens; Qwen3-4B 1.43x and 1.76x (vault lab-prefill-overlap (18)-(26)).
% SOURCE: append, steady band 1.51-1.61x: Llama-3.2-1B at L = 28,672 (1.51x, 1.57x) and Qwen3-4B at L = 16,384 and 28,672 (1.52-1.61x); the 1B's 1.8-2.3x at L = 16,384 includes a slowdown of the existing order at positions ~8k-20k (vault lab-agent-step (2), (4)).
% SOURCE: warm switch S2 <-> S4 7.1 ms (1B) and 7.7 ms (Qwen3-4B), 0.2% of a Qwen3-4B 4k + 4k append on S4 (5,238 ms) (vault lab-split-sweep (5)).
% SOURCE: 16k + 4k append, composed on S4 vs composed on S2: Llama-3.2-1B 5,493.0 -> 2,950.8 ms (1.86x), Qwen3-4B 19,332.6 -> 10,342.7 ms (1.87x); composed on S4 vs the existing host program on S2: Qwen3-4B 31,141.0 -> 10,342.7 ms (3.01x) (vault lab-split-sweep (4)-(5)).
% DECISION: the headline numbers are re-checked against the audit's verdicts at the results freeze (2026-10-21). The Llama-3.2-1B counterpart of the 3.0x (4.27x) appears only in the evaluation, with the existing order's slowdown at positions ~8k-20k named; revisit once that slowdown is isolated (ruling of 2026-10-15).
% DECISION: turn-level speedups (an append plus its decode) are inferred until measured, and the introduction does not state them.

% VALUE: (4) insight: concurrency is free in the hardware but rationed by one in-order CP whose instruction order is the schedule; the model takes per-operator times and a schedule and returns its TTFT, with its accuracy and its split choices
The insight behind these results is that concurrency on a spatial array is free in the hardware but rationed by one in-order command processor.
Commands queued on one hardware context never overlap, whereas one command that issues two groups' work before it waits on either runs them in 1.00--1.06$\times$ the time of the longer one.
The CP's instruction order is therefore the schedule, and its cost can be computed before anything runs.
We capture it in an event model that takes each operator's device time, measured once per configuration, and a candidate schedule, and returns the schedule's TTFT.
The model is within 3.5\% of the measured TTFT on schedules and on a model family it was not fitted on, and it picks the faster split in all 16 configurations measured.
% SOURCE: pair overlap 0.998-1.058 over 9 pairs (vault lab-prefill-overlap (3)); event model within +3.5% on every schedule not used for fitting on both models, -0.3% on Qwen3-4B at 2,048 tokens (vault ch3-formulation, 4k model fixes); split winners 16 of 16 with a least-squares attention line (vault lab-split-sweep (7)).

% VALUE: (5) contributions: 4 noun-phrase items, no numbers; engineering sentence after
In summary, the primary contributions of this work are as follows:
\begin{itemize}
    \item A characterization of where expert designs lose end-to-end LLM time on a spatial NPU, in both prefill and decode.
    \item A measured model of how one in-order command processor turns instruction streams into concurrency on disjoint column groups, with an event model that prices a schedule before it runs.
    \item A formulation of composing concurrent column groups as scheduling through a single in-order server, with a lower bound and the legality conditions that a composed command must meet.
    \item The design and implementation of a composer that makes the instruction-order and spatial-split decisions automatically for unmodified designs, with bit-identical outputs.
\end{itemize}
The composer replaces only the dispatch loop of the design's host program, builds every split from the design's unchanged source, and runs on the existing compiler, driver, and firmware.
